% This is samplepaper.tex, a sample chapter demonstrating the
% LLNCS macro package for Springer Computer Science proceedings;
% Version 2.21 of 2022/01/12
%
\documentclass[runningheads]{llncs}
%
\usepackage[T1]{fontenc}
% T1 fonts will be used to generate the final print and online PDFs,
% so please use T1 fonts in your manuscript whenever possible.
% Other font encondings may result in incorrect characters.
%
\usepackage{graphicx}
% Used for displaying a sample figure. If possible, figure files should
% be included in EPS format.
%
\usepackage{amsmath}
\usepackage{amssymb}
\usepackage{booktabs}
\usepackage{tabularx}
%
% If you use the hyperref package, please uncomment the following two lines
% to display URLs in blue roman font according to Springer's eBook style:
%\usepackage{color}
%\renewcommand\UrlFont{\color{blue}\rmfamily}
%\urlstyle{rm}
%
\begin{document}
%
\title{AI Agentic AI-based Vietnamese Financial News Analysis and Summarization System}
%
%\titlerunning{Abbreviated paper title}
% If the paper title is too long for the running head, you can set
% an abbreviated paper title here
%
\author{Nguyen Vinh Khang\inst{1}\orcidID{0009-0002-4688-0819} \and
Le Quoc Huy\inst{1}\orcidID{0009-0004-4393-1056}}
%
\authorrunning{N. V. Khang et al.}
% First names are abbreviated in the running head.
% If there are more than two authors, 'et al.' is used.
%
\institute{
  University of Science, Ho Chi Minh City National University,
  227 Nguyen Van Cu Street, District 5, Ho Chi Minh City 700000, Vietnam\\
  \email{22125035@student.hcmus.edu.vn}\\
  \email{22125031@student.hcmus.edu.vn}\\
  \url{https://hcmus.edu.vn}
}
%
\maketitle              % typeset the header of the contribution
%
\begin{abstract}
% TODO: Viết sau cùng (sau khi hoàn tất các phần khác).
The abstract should briefly summarize the contents of the paper in
150--250 words.

\keywords{First keyword \and Second keyword \and Another keyword.}
\end{abstract}
%
%
%
% ============================================================
% TODO Section 1 — Introduction (viết sau)
% ============================================================
\section{Introduction}
\label{sec:introduction}

Vietnam's equity market has expanded into one of the most retail-driven markets
in the region. By the end of 2025 domestic investors held roughly $11.9$ million
securities accounts---about $11.6\%$ of the population---after a record year in
which some $2.5$ million new accounts were opened~\cite{vsd2025}, and individual
investors account for the overwhelming majority of participants and for around
$90\%$ of daily trading value~\cite{vinacapital2024}. The market's growing
international standing was confirmed when FTSE Russell announced its
reclassification from Frontier to Secondary Emerging status~\cite{ftse2025}. This
combination---a large, fast-growing base of non-professional investors and rising
institutional interest---creates strong demand for tools that help ordinary
investors make sound decisions amid a flood of fragmented information.

Yet the analytical workflow such investors face remains badly fragmented. The
three classical lenses on a stock---technical signals from price action,
fundamentals from financial statements, and sentiment from the news flow---live
in separate tools with separate interfaces, leaving the investor to reconcile
them manually. The problem is sharper in Vietnamese: relevant company news is
scattered across local outlets and is not readily digested by mainstream,
English-centric tools. Existing retail platforms and robo-advisory services offer
little integrated, Vietnamese-language support that fuses these signals into a
single, horizon-aware recommendation.

A natural reaction is to ``ask a large language model'', but doing so naively is
unsafe in a financial setting. LLMs excel at language understanding and
synthesis, yet they are prone to hallucination, are difficult to make
reproducible, and rarely expose \emph{why} they reached a conclusion---properties
that are unacceptable when the output influences real investment decisions. What
is needed is a hybrid design in which the model is confined to what it does
well---reading and summarizing heterogeneous evidence---while every
decision-critical quantity is produced by explicit, auditable rules.

We present such a system for the Vietnamese market. Its analytical core is a
multi-agent pipeline built on LangGraph: an orchestrator interprets the request
and the investor's risk profile and dispatches, in parallel, a technical, a
fundamental and a news-sentiment agent, whose outputs an aggregator fuses into a
structured buy/wait recommendation with risk-managed price targets. The design is
governed throughout by a single principle---\emph{the language model analyzes,
while deterministic code decides}---which keeps the recommendations explainable
and reproducible. The same data and reasoning back a conversational assistant
that answers data questions by generating sanitized SQL over the live database
and that delegates comprehensive requests to the full analysis pipeline. The
underlying data is kept current by an offline ingestion subsystem, and the
recommendations are validated through a statistically grounded backtesting
framework.

This paper makes the following contributions:
\begin{itemize}
\item A multi-agent architecture, built on LangGraph, that unifies technical,
fundamental and news-sentiment analysis with recommendation generation for
Vietnamese equities within a single, modular pipeline.
\item A fusion mechanism that combines the three sources into a confidence score
by renormalizing weights over only the active sources, paired with a hybrid
LLM/rule decision procedure that keeps verdicts explainable and reproducible.
\item A risk-aware position-sizing scheme that enforces horizon-dependent hard
boundaries on take-profit, stop-loss and holding period, together with a minimum
reward-to-risk constraint.
\item A grounded conversational assistant that answers data questions through
safe text-to-SQL over real market data and escalates holistic requests to the
analysis pipeline.
\item An evaluation methodology based on historical backtesting with statistical
significance testing.
\end{itemize}

\paragraph{Scope.}
We delimit the scope of this work explicitly. The system is a
\emph{decision-support} tool: it produces analysis and recommendations only and
does \emph{not} place or execute orders, connect to any brokerage, or manage
portfolios---acting on its output is left entirely to the user, and its
recommendations do not constitute financial advice. Its coverage is likewise
restricted to the Vietnamese equity market: all price, fundamental and news data,
as well as the supported tickers and indices, are drawn from Vietnamese sources,
and the system is neither designed nor evaluated for other markets.

The remainder of the paper is organized as follows.
Section~\ref{sec:related} reviews related work;
Section~\ref{sec:architecture} presents the system architecture;
Section~\ref{sec:methodology} details the methodology of each agent and the
fusion logic; Section~\ref{sec:implementation} describes the implementation and
data ingestion; Section~\ref{sec:evaluation} reports the evaluation; and
Sections~\ref{sec:discussion} and~\ref{sec:conclusion} discuss limitations and
conclude.

% ============================================================
% TODO Section 2 — Related Work (viết sau)
% ============================================================
\section{Related Work}
\label{sec:related}

Our system sits at the intersection of four lines of research: quantitative and
machine-learning analysis of equities, financial news sentiment analysis, large
language models and agents for finance, and natural-language grounding over
structured data.

\subsection{Quantitative and machine-learning stock analysis}
A large body of work automates the classical lenses of technical and fundamental
analysis and, more recently, applies machine and deep learning to price
forecasting. A representative direction couples news sentiment with price history
in deep architectures, such as FinBERT-LSTM models that feed sentiment-enriched
features into recurrent networks for next-day price prediction~\cite{finbertlstm2022}.
While effective at their stated objective, these models typically optimize a
single target (price or direction), operate as black boxes, and stop short of
producing an explainable, horizon-aware decision that an investor can act on.

\subsection{Financial news sentiment analysis}
Domain-specific language models have become standard for extracting sentiment
from financial text; FinBERT and its successors map news onto sentiment labels
that improve stock-movement prediction~\cite{finbert2023}. For Vietnamese, this
line is comparatively thin but growing: PhoBERT-based classifiers label the
sentiment of stock-article titles with high accuracy~\cite{phobertstock}, and
empirical studies show that sentiment extracted from Vietnamese news measurably
affects market returns~\cite{vnnews2023}. These approaches yield a sentiment
signal in isolation rather than fusing it with technical and fundamental evidence
into a single recommendation, which is the role of our news-sentiment agent.

\subsection{LLMs and agentic systems in finance}
Finance-specific LLMs such as the proprietary BloombergGPT~\cite{bloomberggpt}
and the open-source FinGPT~\cite{fingpt} have demonstrated strong performance on
financial NLP tasks. Building on these, agentic trading frameworks orchestrate
multiple LLM-driven roles: TradingAgents assigns specialized analyst, trader and
risk-manager agents that debate before acting~\cite{tradingagents}, and FinMem
augments an LLM trading agent with layered memory to improve backtested returns
while mitigating hallucination~\cite{finmem}. These frameworks are powerful but
are predominantly built and evaluated on English-language, U.S.-centric data, and
they generally let the language model make the final trading decision---an
arrangement that complicates reproducibility and explainability. Our work differs
by confining the LLM to analysis and delegating every decisive computation to
deterministic rules.

\subsection{Vietnamese-market prediction and decision support}
Research targeting the Vietnamese market has concentrated on index and price
forecasting. Deep architectures including LSTM, bidirectional LSTM and
Transformer variants have been compared for VN-Index prediction~\cite{vndl2025,vnlstm2021},
and conventional machine-learning models have been used to classify price
trends~\cite{vnml2024}. This work is valuable but largely single-source and
predictive: it forecasts a number rather than fusing technical, fundamental and
news signals into an actionable, risk-calibrated recommendation, and it is rarely
LLM-based or explanation-oriented.

\subsection{Grounding LLMs on structured data via text-to-SQL}
Translating natural-language questions into executable SQL is a mature task with
established benchmarks and, increasingly, LLM-based solutions~\cite{text2sqlsurvey,spider2}.
We adopt this idea to ground our conversational assistant: rather than answering
from parametric memory, it generates a sanitized, read-only query against the
live database, so its answers are anchored in real, current data.

\subsection{Positioning}
To the best of our knowledge, no prior system simultaneously (i)~unifies
technical, fundamental and news-sentiment analysis with recommendation generation
for the Vietnamese market, (ii)~adopts a hybrid LLM/rule design that keeps the
decision explainable and reproducible, (iii)~enforces horizon-dependent
risk constraints, and (iv)~validates its recommendations through statistically
grounded backtesting. This paper addresses that gap.

% ============================================================
% TODO Section 3 — System Architecture (viết sau)
% ============================================================
\section{System Architecture}
\label{sec:architecture}

The platform is organized as an end-to-end system that connects a mobile client
to a backend service, an agentic reasoning layer, and a set of data stores and
external providers. A mobile application (built with React Native and Expo)
issues authenticated requests to a FastAPI backend; the backend exposes
conventional market-data and account endpoints and, for analysis, hands control
to an agentic layer implemented on top of LangGraph. The agentic layer reads from
data stores that are kept current by periodic, offline ingestion jobs, so that an
analysis request is served from pre-collected data rather than blocking on live
crawling. Figure~\ref{fig:architecture} shows this organization, which the rest
of this section describes layer by layer.

\begin{figure}[t]
\centering
\includegraphics[width=\textwidth]{fig1_architecture}
\caption{High-level system architecture. Solid online paths serve analysis
requests; the dashed offline path periodically ingests and enriches market,
fundamental and news data into the stores that the agents read from.}
\label{fig:architecture}
\end{figure}

\subsection{Layered backend design}
\label{sec:layered}
The backend follows a strict three-tier separation: \emph{routes} expose REST
endpoints, \emph{services} hold the business logic, and a \emph{data} tier wraps
the database and external APIs. Requests enter through a set of routers---among
them authentication, market data, fundamental analysis, technical indicators,
articles, company profiles, risk appetite, portfolio, and the agentic
endpoints---each delegating to a matching service so that transport concerns stay
out of the domain logic. Protected routes are guarded by JWT-based authentication
enforced in middleware, and every endpoint returns a single standardized response
envelope (a data payload plus an error code, description, request identifier and
success flag), which keeps the client's handling uniform across the whole API.
This conventional outer shell is what hosts the agentic layer described next.

\subsection{The agentic layer}
\label{sec:agentic-layer}
Analytical intelligence is concentrated in two LangGraph \texttt{StateGraph}s that
share the same LLM service (OpenAI's \texttt{GPT-4o}) and data tier. The first is
the \emph{analysis pipeline}: an orchestrator parses the request and fans out, in
parallel, to a technical, a fundamental and a news agent, whose outputs are
merged and synthesized by an aggregator into a structured recommendation. The
agents communicate through a single typed state object, and the parallel results
are combined with a merge reducer so that no agent overwrites another's
contribution. The second graph powers the \emph{conversational assistant}: an
intent classifier routes each message to a chat, knowledge-QA, or market
(text-to-SQL) agent, or delegates a comprehensive request to the analysis
pipeline above; conversation history is persisted so that context survives across
turns. Both graphs deliberately confine the LLM to language understanding and
synthesis, while decision-critical computation is performed in deterministic
code. Section~\ref{sec:methodology} details the individual agents.

\subsection{Data sources and ingestion}
\label{sec:arch-data}
The agents never call external providers on the critical path; they read from
stores that are populated ahead of time. Daily and intraday price data originate
from the SSI FC Data API. Company fundamentals, news articles, and reference
metadata are held in Supabase (managed PostgreSQL), which also stores the
conversational checkpoints; a Redis cache holds short-lived session and token
state. Keeping the data fresh is the job of an \emph{offline ingestion subsystem}:
scheduled jobs crawl Vietnamese financial-news sources and pull price and
fundamental updates, and---for news---an LLM enrichment step extracts the
mentioned tickers, a sentiment label and a short summary before the records are
written and linked to the relevant stocks (Section~\ref{sec:news}). Separating
this offline path from the online analysis path means recommendations are
computed over consistent, pre-validated data and are not delayed by network
latency or source outages. The concrete ingestion mechanism---schedules, crawling
and enrichment---is detailed in Section~\ref{sec:implementation}.

\subsection{Request lifecycle}
\label{sec:lifecycle}
A typical ``analyze ticker $X$'' request illustrates how the layers cooperate.
The client sends an authenticated call carrying the ticker, the user's risk
profile and any data-source selections. The backend validates the request and
invokes the analysis graph: the orchestrator derives the analysis window and
candle interval, the three agents read the relevant slices of the pre-ingested
data and produce their partial results in parallel, and the aggregator fuses them
into a confidence-scored buy/wait recommendation with risk-managed price targets.
The structured result is returned through the standard response envelope and
rendered by the mobile client. The online request thus touches only the stores,
while all crawling and enrichment happen earlier on the offline path.: Figure 1 (LangGraph StateGraph), AgentState, data sources.

% ============================================================
% Section 4 — Methodology
% ============================================================
%
\section{Methodology}
\label{sec:methodology}

We frame stock decision support as a \emph{multi-source fusion} problem:
given a target ticker and an investor risk profile, the system produces a
structured, actionable recommendation that integrates technical, fundamental
and news-sentiment evidence. The analysis pipeline is implemented as a
LangGraph \texttt{StateGraph} of five cooperating agents---an
\emph{orchestrator}, three parallel analysis agents (technical, fundamental,
news), and an \emph{aggregator}---that share a single typed state object.
A guiding design principle runs through every stage: \emph{the language model
analyzes, while deterministic Python code decides}. Large language models
(LLMs) are used where natural-language understanding and synthesis are
required (request parsing, news summarization, narrative justification),
whereas all decision-critical quantities---indicator scores, fusion
confidence, buy/wait verdicts and risk boundaries---are computed by explicit,
auditable rules. This separation keeps the system's outputs reproducible and
explainable, two properties that are essential in a financial advisory
setting. The remainder of this section details each agent in turn.

\subsection{Orchestrator Agent}
\label{sec:orchestrator}

The orchestrator is the entry node of the pipeline. Its responsibility is to
translate a free-form user request together with the investor's risk profile
into a concrete, machine-readable \emph{plan} that parameterizes the
downstream analysis agents. Rather than letting the LLM emit unconstrained
text, we use \emph{structured decoding}: the model is asked to populate a
Pydantic schema, so its output is guaranteed to be a valid plan object. The
plan specifies, for each agent, the time window to analyze and---for the
technical agent---the candle interval. We invoke \texttt{GPT-4o} with
temperature $0$ to make the planning step as deterministic as possible.

\paragraph{Shared state.}
All five agents communicate through a single typed state object,
\texttt{AgentState}, defined as a \texttt{TypedDict}. It carries the execution
\texttt{mode} (\texttt{auto}, \texttt{manual} or \texttt{chat}), the raw
\texttt{user\_input}, the \texttt{risk\_appetite}, the target \texttt{symbol},
a \texttt{data\_selection} object (which lets the user toggle individual data
sources and indicators), the orchestrator-generated \texttt{plan}, the
collected \texttt{agent\_results}, and the \texttt{final\_output}. The results
field is annotated with the \texttt{operator.or\_} reducer so that the three
analysis agents, which run concurrently, can each write their partial result
into the shared dictionary without overwriting one another. The orchestrator
fans out to the technical, fundamental and news agents in parallel; once all
three complete, their merged results flow into the aggregator
(see Section~\ref{sec:architecture}, Fig.~\ref{fig:architecture}).

\paragraph{Risk-horizon mapping.}
The investor's risk profile reaches the system as a \texttt{period} field that
takes one of three canonical values---\texttt{short\_term}, \texttt{mid\_term}
or \texttt{long\_term}---which map directly onto the horizon
$h \in \{\textsf{short}, \textsf{mid}, \textsf{long}\}$. When the period is
missing or unrecognized, the system falls back to \textsf{short}, the most
conservative horizon. The horizon then determines two quantities: the
look-back window $[\,t_{\text{from}}, t_{\text{to}}\,]$ used to gather
evidence, and the candle interval used for technical analysis. We always set
$t_{\text{to}}$ to the current date and compute $t_{\text{from}}$ by
subtracting the horizon-specific look-back length.
Table~\ref{tab:horizon} summarizes this mapping.

\paragraph{Rationale for the interval and look-back choices.}
The two parameters are chosen jointly so that the technical analysis operates
on the same time scale as the investor's intended holding period, and so that
each indicator receives enough history to be statistically meaningful without
diluting the signal with stale data. The guiding principle is that the candle
interval should resolve the price movements relevant to the horizon while
suppressing higher-frequency fluctuations that act as noise at that horizon. A
short-term trader cares about day-to-day swings, so daily (\texttt{1d}) candles
preserve exactly the granularity that drives short-horizon entries and exits; a
60-day window keeps the analysis focused on the prevailing short-term regime
and avoids letting older price action distort fast-moving indicators such as
RSI or the MACD histogram. For a multi-month horizon, daily candles would
introduce noise that is irrelevant to a swing position, so we aggregate to
weekly (\texttt{1w}) candles, which expose the medium-term trend; the
180--365-day window then supplies enough weekly observations to warm up the
slower moving averages and oscillators reliably. For a long-term, position-style
horizon we move to monthly (\texttt{1M}) candles, which strip out intra-year
fluctuations and reveal the primary trend, and we require at least two years
($\geq 730$ days) of history so that the long-run pattern is established on a
sufficient number of monthly observations. In every case the look-back length
is deliberately bounded: a longer window would blend together distinct market
regimes and weaken the relevance of the resulting signal to current conditions.

\begin{table}[t]
\centering
\caption{Mapping from the normalized investment horizon to the technical
candle interval and the evidence look-back window. The candle interval is
enforced deterministically in code rather than trusted to the LLM.}
\label{tab:horizon}
\begin{tabular}{@{}llll@{}}
\toprule
Horizon $h$ & Description & Candle interval & Look-back window \\
\midrule
\textsf{short} & days--weeks      & 1d & last 60 days \\
\textsf{mid}   & several months   & 1w & last 180--365 days \\
\textsf{long}  & $>1$ year        & 1M & last $\geq 730$ days \\
\bottomrule
\end{tabular}
\end{table}

\paragraph{Deterministic interval enforcement.}
Although the LLM proposes an interval as part of its structured output, we do
not rely on it for the final value. After decoding the plan, the candle
interval is overwritten in Python according to the deterministic map
\[
\textsf{interval}(h) =
\begin{cases}
\texttt{1d} & h = \textsf{short},\\
\texttt{1w} & h = \textsf{mid},\\
\texttt{1M} & h = \textsf{long}.
\end{cases}
\]
This guarantees that each horizon is genuinely analyzed on a different candle
granularity, independently of any drift in the model's behavior. The
underlying market database stores only daily (\texttt{1d}) candles; the market
service aggregates these on demand into weekly (\texttt{1w}) and monthly
(\texttt{1M}) candles. The orchestrator also enforces two further invariants
that are critical for downstream correctness: it never invents a ticker symbol
that is not present in the input, and it always emits parameters for all three
analysis agents so that the fan-out is complete. In \texttt{manual} mode, where
an advanced user supplies an explicit plan, the orchestrator bypasses the LLM
entirely and forwards the user-provided plan unchanged.

\subsection{Technical Analysis Agent}
\label{sec:technical}

The technical agent quantifies the price action of the target ticker over the
planned window. It computes five widely used indicators on the candle series
delivered by the market service, converts each indicator into a binary bullish
signal through a transparent rule set, and aggregates these into an integer
score. Crucially, the output is not a single opaque number: for every indicator
the agent emits its raw values together with a short natural-language
justification, so that the downstream aggregator---and ultimately the user---can
trace exactly why a signal fired. All indicators are implemented in pure
NumPy/Pandas rather than through a third-party technical-analysis library, which
lets us calibrate the smoothing conventions to match the values displayed by the
domestic data provider and popular charting platforms.

\subsubsection{Indicator computation.}
Let $c_t$, $h_t$ and $l_t$ denote the close, high and low of candle $t$. The
five indicators are defined as follows.

\paragraph{Relative Strength Index (RSI, period $n=14$).}
With $\Delta_t = c_t - c_{t-1}$, the average gain and loss are seeded by the
simple mean of the first $n$ deltas and then updated with Wilder's smoothing:
\begin{equation}
\overline{G}_t = \frac{(n-1)\,\overline{G}_{t-1} + G_t}{n}, \qquad
\overline{L}_t = \frac{(n-1)\,\overline{L}_{t-1} + L_t}{n},
\end{equation}
where $G_t = \max(\Delta_t, 0)$ and $L_t = \max(-\Delta_t, 0)$. The index is
\begin{equation}
\mathrm{RSI}_t = 100 - \frac{100}{1 + \overline{G}_t / \overline{L}_t}.
\end{equation}

\paragraph{Moving averages (SMA and EMA).}
The $p$-period simple moving average is
$\mathrm{SMA}^p_t = \frac{1}{p}\sum_{i=0}^{p-1} c_{t-i}$; we use $p\in\{20,50\}$.
The exponential moving average is seeded with $\mathrm{SMA}^p$ and updated with
$k = 2/(p+1)$:
\begin{equation}
\mathrm{EMA}^p_t = k\,c_t + (1-k)\,\mathrm{EMA}^p_{t-1}.
\end{equation}

\paragraph{Moving Average Convergence Divergence (MACD, $12/26/9$).}
\begin{equation}
\mathrm{MACD}_t = \mathrm{EMA}^{12}_t - \mathrm{EMA}^{26}_t, \quad
\mathrm{Sig}_t = \mathrm{EMA}^{9}(\mathrm{MACD})_t, \quad
\mathrm{Hist}_t = \mathrm{MACD}_t - \mathrm{Sig}_t.
\end{equation}

\paragraph{Bollinger Bands (period $20$, $2\sigma$).}
The middle band is $\mathrm{SMA}^{20}_t$ and the envelopes use the population
standard deviation $\sigma_t$ over the same window (matching the convention of
the data provider and TradingView):
\begin{equation}
\mathrm{BB}^{\mathrm{mid}}_t = \mathrm{SMA}^{20}_t, \qquad
\mathrm{BB}^{\mathrm{up/lo}}_t = \mathrm{SMA}^{20}_t \pm 2\sigma_t.
\end{equation}

\paragraph{Stochastic KDJ ($9/3/3$).}
With the raw stochastic value over the last $9$ candles
\begin{equation}
\mathrm{RSV}_t = \frac{c_t - \min_{9} l}{\max_{9} h - \min_{9} l}\times 100,
\end{equation}
the smoothed lines (Wilder factor $\tfrac{1}{3}$, seeded at $50$) are
\begin{equation}
K_t = \tfrac{1}{3}\mathrm{RSV}_t + \tfrac{2}{3}K_{t-1}, \quad
D_t = \tfrac{1}{3}K_t + \tfrac{2}{3}D_{t-1}, \quad
J_t = 3K_t - 2D_t.
\end{equation}

\subsubsection{Rule-based scoring.}
Each enabled indicator $i$ is mapped to a binary score $s_i\in\{0,1\}$ by a small
set of bullish conditions; $s_i = 1$ whenever \emph{any} of the indicator's
conditions hold, evaluated by comparing the most recent candle against the
previous one. Table~\ref{tab:scoring} lists these conditions. The design favours
\emph{nascent} upward momentum---fresh crossovers, accelerating gaps, and
rebounds from oversold zones---rather than already-overextended readings, which
are explicitly treated as non-bullish (e.g.\ $\mathrm{RSI}\geq 70$ or price above
the upper Bollinger band).

\begin{table}[t]
\centering
\caption{Bullish scoring rules. An indicator earns $s_i=1$ if any listed
condition holds between the previous candle and the current one.}
\label{tab:scoring}
\begin{tabularx}{\textwidth}{@{}lX@{}}
\toprule
Indicator & Conditions for $s_i = 1$ (logical OR) \\
\midrule
RSI  & crosses above $50$; or $50\!\leq\!\mathrm{RSI}\!\leq\!70$ and rising; or $\mathrm{RSI}\!<\!35$ and rising \\
MA   & golden cross (SMA20 crosses above SMA50); or SMA20$>$SMA50 with widening gap; or price$>$SMA20$>$SMA50 \\
BOLL & price crosses above middle band; or above middle with increasing distance; or rebound off lower band \\
MACD & crosses above signal; or MACD$>$signal with rising histogram; or histogram turns positive \\
KDJ  & $K$ crosses above $D$; or $K>D$ with rising $J$; or $K<20$ and rising \\
\bottomrule
\end{tabularx}
\end{table}

\subsubsection{Aggregation and output.}
The agent reports the total score together with the maximum attainable score,
\begin{equation}
S_{\mathrm{tech}} = \sum_{i \in \mathcal{E}} s_i, \qquad
S_{\max} = |\mathcal{E}|,
\end{equation}
where $\mathcal{E}$ is the set of indicators the user has enabled through the
\texttt{data\_selection} object (all five by default). Reporting $S_{\max}$
alongside $S_{\mathrm{tech}}$ lets the aggregator normalize the technical signal
to $[0,1]$ regardless of how many indicators are active
(Section~\ref{sec:aggregator}). The agent finally packages the symbol, interval,
analyzed period, the latest price, $S_{\mathrm{tech}}/S_{\max}$, and the
per-indicator values and reasons into the shared state. In live mode the latest
price is taken from the most recent one-minute candle to reflect the true
current quote, while the indicators themselves are computed on the
horizon-appropriate candles.

\subsection{Fundamental Analysis Agent}
\label{sec:fundamental}

While the technical agent reads price action, the fundamental agent assesses the
underlying financial health of the issuer. It queries the company's annual
financial records and condenses them into a structured report covering five
groups of metrics. A deliberate design choice distinguishes this agent from the
technical one: its evidence window is \emph{independent of the investor's
horizon}. Whereas price-based signals must be re-scaled to the trading horizon,
fundamentals evolve slowly---they are published only quarterly or annually and
reflect the firm's structural condition---so the agent always retrieves the most
recent \emph{five fiscal years} of fundamentals for the short-, mid- and
long-term profiles alike. Five years is long enough to span a full business
cycle and to expose multi-year trends in profitability and growth, yet short
enough to remain representative of the company in its current form.

\paragraph{Data source and metric groups.}
The agent reads per-ticker financial data from the platform's database, drawing
on dedicated tables for the analyst summary, valuation/profitability indicators,
the income statement, the cash-flow statement and the balance sheet. From these
it assembles five metric groups, each of which the user can enable or disable
through the same \texttt{data\_selection} mechanism used by the technical agent.
Table~\ref{tab:fundamental} lists the groups and their constituent metrics. The
agent presents the most recent fiscal year as the headline snapshot---together
with year-over-year growth---while the five-year history provides the trend
context that informs the qualitative summary.

\begin{table}[t]
\centering
\caption{The five fundamental metric groups and their constituent indicators.
Each group can be toggled independently by the user.}
\label{tab:fundamental}
\begin{tabularx}{\textwidth}{@{}lX@{}}
\toprule
Group & Metrics \\
\midrule
Valuation         & P/E, P/B, EV/EBITDA, EPS \\
Profitability     & ROE, ROA, gross margin, net margin \\
Growth            & revenue YoY, profit YoY, net revenue, net profit \\
Financial health  & current ratio, debt-to-equity, interest coverage \\
Cash flow         & operating cash flow (CFO), capital expenditure (CAPEX), dividends paid \\
\bottomrule
\end{tabularx}
\end{table}

\paragraph{Sector awareness and robustness.}
Several ratios are ill-defined for particular sectors---gross margin, current
ratio and debt-to-equity, for instance, do not carry their usual meaning for
banks---so the agent detects such cases and annotates the report accordingly
rather than emitting misleading figures. Robustness to missing data is treated
as a first-class concern: newly listed tickers or companies without filed
reports may lack some or all fundamentals. A data-layer error for a single
ticker must never abort the pipeline, which is essential when the system is run
across an entire index such as VN30 or VN100. When fundamentals are missing or a
query fails, the agent emits an explicit ``no data'' marker instead of raising
an exception; the aggregator recognizes this marker and simply drops the
fundamental source from the confidence computation (Section~\ref{sec:aggregator})
rather than penalizing the stock.

\paragraph{Output.}
The agent's output is a structured, human-readable report---the headline year,
the analyst summary, and the enabled metric groups. This narrative is consumed by
the aggregator, which interprets it into a discrete health label
$\textsf{fundamental\_health}\in\{\textsf{strong},\textsf{neutral},\textsf{weak},
\textsf{N/A}\}$ used in the final fusion. Keeping the classification in the
aggregator, rather than hard-coding thresholds here, lets the system weigh the
fundamental picture jointly with the technical and news evidence.

% ============================================================
\subsection{News Sentiment Agent}
\label{sec:news}

The third analysis agent captures market-relevant information that is invisible
to price and accounting data alike: company news. It follows a two-stage design
that separates expensive, reusable preprocessing from lightweight,
query-specific synthesis.

\paragraph{Stage 1: offline ingestion and enrichment.}
Financial news articles concerning Vietnamese listed companies are collected
into the platform's database, where each record stores the title, link,
publication time, source, and full text. During ingestion every article is
passed through an LLM enrichment step that populates a structured schema,
\texttt{NewsExtraction}, consisting of (i) the list of stock tickers mentioned
in the article, (ii) a sentiment label
$\in\{\textsf{positive},\textsf{neutral},\textsf{negative}\}$, and (iii) a
concise one-to-three-sentence summary. The extracted tickers drive a
many-to-many association between articles and stocks (an
\texttt{Article\_Stock} link table), so that one article may be attached to
several companies. We distinguish \emph{exclusive} articles, which mention a
single ticker and are therefore used for company-specific analysis, from
\emph{macro} articles, which reference two or more tickers and feed broader
market context. Performing extraction once at ingestion time means the
per-article sentiment and summary are cached and reused across all subsequent
analyses, rather than being recomputed on every request.

\paragraph{Stage 2: online synthesis.}
When an analysis is requested for a ticker, the agent retrieves the articles
linked to that ticker and filters them to the date window
$[\,t_{\text{from}}, t_{\text{to}}\,]$ supplied by the orchestrator. The
filtered set---each item carrying its date, title, cached summary and cached
sentiment---is presented to \texttt{GPT-4o}, which is instructed to first
discard items that are irrelevant or uninformative and then synthesize the
remainder into a structured brief comprising (1)~a summary of the principal
news, (2)~the positive points, and (3)~the negative points or risks. If no
useful article remains, the agent returns an explicit ``no useful news''
response. This synthesis step is the one place in the pipeline where we permit a
small amount of generative freedom (decoding temperature $0.2$) so that the
brief reads fluently; it nonetheless remains an \emph{analysis} step, since the
final decision is still taken downstream by deterministic rules.

\paragraph{Output and robustness.}
As with the other agents, the news source can be toggled off by the user, and
the absence of articles in the requested window is handled gracefully rather
than as an error. The textual brief is consumed by the aggregator, which
condenses it into a single discrete signal
$\textsf{article\_sentiment}\in\{\textsf{positive},\textsf{neutral},
\textsf{negative},\textsf{N/A}\}$. When news is disabled or unavailable, the
\textsf{N/A} value causes the source to be dropped from the confidence
computation (Section~\ref{sec:aggregator}) without penalizing the stock.

\subsection{Aggregator and Decision Logic}
\label{sec:aggregator}

The aggregator is the decision node where the three analysis streams are fused
into a single, actionable recommendation. It is also where the design principle
stated at the start of this section is enforced most strictly. An LLM
(\texttt{GPT-4o}) is used to read the heterogeneous agent outputs and emit a
structured \texttt{InvestmentRecommendation} object, including a Vietnamese
narrative justification. However, every decision-critical quantity---the fusion
confidence, the final buy/wait verdict, and the risk-managed price targets---is
recomputed and, where necessary, overridden by deterministic Python rules that
do not depend on the model's discretion. The LLM is given the same rules in its
prompt, but the Python layer is authoritative: this defence-in-depth keeps the
system reproducible and prevents an over-optimistic model from issuing an unsafe
recommendation.

\paragraph{Signals and data-presence flags.}
The aggregator works from three raw signals: the integer technical score
$S_{\mathrm{tech}}\in[0,S_{\max}]$ read \emph{verbatim} from the technical agent
(never recomputed), a fundamental health label
$f\in\{\textsf{strong},\textsf{neutral},\textsf{weak},\textsf{N/A}\}$, and a news
sentiment label $a\in\{\textsf{positive},\textsf{neutral},\textsf{negative},
\textsf{N/A}\}$, both inferred by the LLM from the respective agents' reports.
Crucially, whether a source actually carries data is decided in Python, not by
the model: boolean flags $\delta_{\mathrm{tech}},\delta_{\mathrm{fund}},
\delta_{\mathrm{news}}$ are set by inspecting the technical indicators and by
scanning the fundamental and news texts for explicit ``no data'' markers. This
prevents the common failure mode in which a model lazily declares data missing
when it is in fact present.

\subsubsection{Confidence fusion.}
Each source is first normalized to a component in $[0,1]$:
\begin{equation}
c_{\mathrm{tech}} =
\begin{cases}
S_{\mathrm{tech}}/S_{\max} & S_{\max} > 0,\\
0.5 & S_{\max} = 0,
\end{cases}
\qquad
c_{\mathrm{fund}} = \phi_f(f), \qquad
c_{\mathrm{news}} = \phi_a(a),
\end{equation}
where the label-to-score maps assign a bullish reading to $1$, a neutral or
missing reading to $0.5$, and a bearish reading to $0$:
\begin{align}
\phi_f:\ & \textsf{strong}\!\mapsto\!1,\ \textsf{neutral}\!\mapsto\!0.5,\
\textsf{weak}\!\mapsto\!0,\ \textsf{N/A}\!\mapsto\!0.5,\\
\phi_a:\ & \textsf{positive}\!\mapsto\!1,\ \textsf{neutral}\!\mapsto\!0.5,\
\textsf{negative}\!\mapsto\!0,\ \textsf{N/A}\!\mapsto\!0.5.
\end{align}
The three sources carry equal base weights
$w_{\mathrm{tech}}=w_{\mathrm{fund}}=w_{\mathrm{news}}=\tfrac{1}{3}$. The central
idea is that weights are \emph{renormalized over the active sources only}. Let
$\mathcal{A}=\{i : \delta_i = \text{true}\}$ be the set of sources that are
enabled and carry data. The effective weights and the confidence are
\begin{equation}
\tilde{w}_i = \frac{w_i}{\sum_{j\in\mathcal{A}} w_j}\ \text{for } i\in\mathcal{A}
\ (\text{else } 0), \qquad
\mathrm{conf} = \sum_{i\in\mathcal{A}} \tilde{w}_i\, c_i \in [0,1].
\label{eq:confidence}
\end{equation}
Renormalization is what makes the score robust to missing or disabled sources:
rather than letting an absent source contribute a neutral $0.5$ that drags the
result toward the midpoint, its weight is redistributed across the surviving
sources, so each active source spans the full $[0,1]$ range. With uniform base
weights this reduces to the mean of the active components,
$\mathrm{conf} = \frac{1}{|\mathcal{A}|}\sum_{i\in\mathcal{A}} c_i$; when only a
single source is active the confidence equals that source's own component.

\subsubsection{Buy/wait decision.}
The verdict combines a mandatory technical \emph{trigger} with two confirmatory
\emph{filters}. The technical gate requires the score to reach $60\%$ of the
attainable maximum,
\begin{equation}
\tau_{\mathrm{tech}} = \lceil 0.6\, S_{\max} \rceil, \qquad
g_{\mathrm{tech}} = \big(S_{\max}=0\big)\ \lor\ \big(S_{\mathrm{tech}} \geq
\tau_{\mathrm{tech}}\big),
\end{equation}
where the disjunct $S_{\max}=0$ disables the gate if the user has turned off all
technical indicators. The fundamental and news streams act purely as vetoes: a
position is blocked only by an \emph{explicitly} adverse reading, while neutral
and N/A values never block. The decision is therefore
\begin{equation}
\textsf{Buy} \iff
g_{\mathrm{tech}} \ \land\
\neg(\delta_{\mathrm{fund}} \land f{=}\textsf{weak}) \ \land\
\neg(\delta_{\mathrm{news}} \land a{=}\textsf{negative}) \ \land\
\mathrm{conf} \geq \tau_{\mathrm{conf}},
\end{equation}
with the confidence threshold $\tau_{\mathrm{conf}} = 0.55$; any other case
yields \textsf{Wait}. This asymmetry---a strong technical trigger to enter, but
only clear negatives to abstain---reflects that missing evidence should make the
system cautious rather than punitive. All of these conditions are evaluated in
Python after the LLM responds and override the model's own verdict when they
disagree.

\subsubsection{Risk-managed price targets.}
For a \textsf{Buy}, the entry price is anchored to the latest quote, and the
take-profit (TP), stop-loss (SL) and maximum holding period are constrained to
hard, horizon-dependent boundaries (Table~\ref{tab:bounds}). The LLM may
fine-tune levels using support/resistance or volatility cues, but its proposals
are clamped in Python: each target is converted to a percentage from entry,
clipped to the allowed band, and the price recomputed,
\begin{equation}
\mathrm{TP} = p_e\big(1 + \mathrm{clip}(r_{\mathrm{tp}}, \text{band})\big),
\qquad
\mathrm{SL} = p_e\big(1 - \mathrm{clip}(r_{\mathrm{sl}}, \text{band})\big),
\end{equation}
where $p_e$ is the entry price. The system additionally requires a
reward-to-risk ratio of at least $1.5$; a candidate that cannot satisfy this
within the hard bounds is downgraded to \textsf{Wait}. For a \textsf{Wait}, all
price fields are left null.

\begin{table}[t]
\centering
\caption{Hard risk boundaries by investment horizon, expressed as a percentage
of the entry price for take-profit (TP) and stop-loss (SL), and as a number of
candles for the maximum holding period.}
\label{tab:bounds}
\begin{tabular}{@{}llll@{}}
\toprule
Horizon & TP (\%) & SL (\%) & Max hold (candles) \\
\midrule
\textsf{short} & 8--15  & 4--7   & 5--15 \\
\textsf{mid}   & 15--30 & 7--12  & 15--60 \\
\textsf{long}  & 40--80 & 15--20 & 120--260 \\
\bottomrule
\end{tabular}
\end{table}

\paragraph{Output and robustness.}
The agent returns an \texttt{InvestmentRecommendation} containing the verdict,
the (possibly null) entry/TP/SL prices and holding period, a four-part narrative
(technical, fundamental, news, and a qualitative summary), the raw signals, and
the computed confidence together with its weight breakdown. The exact,
post-clamping price line is injected into the summary so that the narrative can
never contradict the displayed numbers. Two robustness mechanisms protect the
node: if structured decoding is truncated by the output-token cap, the request
is retried once with an instruction to be terse; and any unhandled error yields
a safe default of \textsf{Wait} with null targets rather than a malformed
response.

\subsection{Conversational Assistant}
\label{sec:chatbot}

Alongside the on-demand analysis pipeline, the platform exposes a conversational
assistant for free-form questions. It is implemented as a separate LangGraph
\texttt{StateGraph} whose conversation history is persisted to PostgreSQL through
a checkpointer keyed by thread, so that multi-turn context survives across
requests. The graph begins with an \emph{intent classifier}
(\texttt{GPT-4o}, temperature $0$, structured output) that reads the latest
message together with the last six turns of history and routes it to one of four
branches: \textsf{out-of-scope} questions are politely declined; \textsf{greeting}
and small-talk are handled by a lightweight chat agent; conceptual
\textsf{knowledge} questions (e.g.\ ``what is RSI?'') are answered by a QA agent
without touching the database; and \textsf{market} questions, which require real
data, are sent to the market agent described next.

\paragraph{Grounded question answering via text-to-SQL.}
The market agent answers data questions by querying the database directly rather
than relying on the model's parametric memory, in a four-stage pipeline. The
incoming question is first Unicode-normalized (NFC) so that Vietnamese text typed
with different diacritic compositions tokenizes consistently.

In the \emph{first} stage the agent generates SQL. It is given the live database
schema---the table and column definitions, introspected dynamically as DDL---and
the current date in Vietnam time, and is asked to translate the question into a
single read-only statement, returned as a structured object comprising the SQL
string, a flag for whether technical indicators are required, an optionally
resolved ticker, a reason field used when no query is possible, and the detected
language of the question. The prompt encodes the schema's conventions so that the
generated query is correct by construction: only a single \texttt{SELECT} (or
\texttt{WITH \ldots\ SELECT}) over existing tables and columns, case-sensitive
identifiers wrapped in double quotes, uppercased tickers, a bounded
\texttt{LIMIT}, joins from the fundamental tables to the stock table by foreign
key, and news joined through the \texttt{Article\_Stock} link table. A
domain-specific correctness rule is also enforced in the query itself: the
intraday and daily price tables store OHLCV values pre-divided by $1000$, so the
model is required to multiply these columns by $1000$ inside the SQL (preserving
the column aliases) so that the returned figures are the true prices. If the
model resolves a ticker but nevertheless returns an empty query---a known
failure mode on open-ended ``is this stock worth buying?'' phrasings---a
deterministic fallback substitutes a default query against the company's
fundamental-summary table, so a resolvable ticker always yields data rather than
a refusal.

In the \emph{second} stage the query is executed. It is never run blindly: it
passes through a sanitizer that permits only read-only \texttt{SELECT}s over an
allow-list of tables and runs under a timeout, so the model never touches the
database directly. An unsafe query is refused outright, whereas a syntax or
column error triggers a single guided retry in which the error message is fed
back to the model for correction. \emph{Third}, because indicators such as RSI,
MACD and KDJ are not stored, a question flagged as needing technical computation
fetches the raw daily candles by SQL (ordered most-recent-first, enough bars to
satisfy the fifty-bar minimum) and computes the indicators on the fly, reusing
the service of Section~\ref{sec:technical} and keeping only the latest value of
each. \emph{Fourth}, the retrieved rows---and any computed indicators---are handed
to a second \texttt{GPT-4o} call that composes a natural-language answer
\emph{grounded entirely in the data}, in the user's own language, with explicit
instructions not to fabricate figures and to state the date of any time-stamped
data. As elsewhere, the boundary between generation and decision is deliberate:
classification and SQL generation run deterministically (temperature $0$), while
only the final answer synthesis is given mild creative latitude.

\paragraph{Delegation to the analysis pipeline.}
A distinct route is reserved for \emph{comprehensive} requests. When the
classifier recognizes that the user is asking for a full, multi-factor evaluation
of a ticker---``should I buy VIC?'', ``give me a complete assessment of HPG''---a
single SQL lookup is insufficient. In such cases the assistant delegates to the
end-to-end analysis pipeline of Sections~\ref{sec:orchestrator}--%
\ref{sec:aggregator}, running the orchestrator, the three analysis agents and the
aggregator, and returns the resulting structured recommendation through the chat
interface. This lets a single conversational entry point serve both quick factual
look-ups and the system's full analytical capability.

% ============================================================

\section{Implementation}
\label{sec:implementation}

This section records the concrete technology choices and the engineering details
that make the system operational, with particular attention to the offline
ingestion subsystem introduced in Section~\ref{sec:arch-data}.

\subsection{Backend stack and deployment}
The backend is a FastAPI~0.128 application served by Uvicorn on Python~3.12, with
all request and response models defined as Pydantic~v2 schemas. The agentic layer
is built on LangGraph~1.0, and conversational state is persisted with the
companion PostgreSQL checkpointer. Three stores back the service: Supabase
(managed PostgreSQL) is accessed both through its REST client and, for the
chatbot checkpoints and text-to-SQL queries, through a pooled \texttt{psycopg}
connection; a Redis instance holds short-lived session and token state. The
service is deployed as a single web process.

\subsection{LLM integration and structured output}
All agents share a thin client over OpenAI's \texttt{GPT-4o}. Wherever an agent
must return machine-consumable data---the orchestrator plan, the intent label,
the generated SQL, the aggregator's recommendation---we use the SDK's structured
decoding, which constrains the model to emit a valid instance of a Pydantic
schema rather than free text. Sampling temperature is chosen per role to match
the earlier ``analyze vs.\ decide'' principle: deterministic steps (orchestration,
intent classification, SQL generation, the aggregator's signal extraction) run at
temperature $0$, while the genuinely generative steps (news synthesis, the
aggregator's narrative, chat) are allowed a small temperature of $0.2$--$0.3$.
The aggregator additionally caps its output tokens and retries once with a
terse-output instruction if structured decoding is truncated. A Google Gemini
client is also configured through \texttt{langchain-google-genai} as an
alternative provider.

\subsection{Market-data integration}
At request time, price data is served by an SSI FC Data API client embedded in
the backend. The client manages bearer-token authentication and transparently
refreshes the token and retries once on an authentication failure. To avoid
redundant synchronization under bursty traffic, it applies a 120-second
per-(symbol, interval) cooldown, and it guards the periodic refresh of the most
active symbols with a lock. Because the database stores only daily candles, the
client aggregates them on demand into the weekly and monthly series required by
the mid- and long-term horizons.

\subsection{Data ingestion subsystem}
\label{sec:ingestion}
Keeping the stores current is the responsibility of two \emph{independent}
ingestion services---one for prices and one for news---each packaged as its own
Docker image and deployed separately. Rather than relying on the operating
system's cron, each service runs a long-lived \texttt{scheduler.py} that polls
the wall-clock time (in Vietnam time, UTC$+7$) and triggers work at the
appropriate moments, which keeps scheduling logic and trading-calendar awareness
in one place.

\paragraph{Price ingestion.}
On trading days (Monday--Friday) between the 09:00 open and the 15:00 close, the
price scheduler supervises three SSI WebSocket streams---one-minute candles,
daily candles, and the market indices---restarting any stream that crashes.
Shortly after the close (15:05) it tears the streams down and runs, once per day,
three \emph{standardization} jobs that reconcile the streamed data into the
canonical price tables and retain a rolling five-year daily history. Separate
\emph{initialization} scripts backfill that five-year history on first setup,
fetching in thirty-day chunks (the SSI per-request limit) with a short delay
between calls to respect rate limits. The tracked universe is the VN30
constituents, resolved from the index membership tables, and the documented
$\times 1000$ price-scaling convention is applied so that stored values match the
provider's representation.

\paragraph{News ingestion.}
The news scheduler runs weekly (Monday, 00:00). For each listed symbol it
discovers candidate articles through a news search API, restricted to a small set
of trusted Vietnamese financial outlets, extracts the article body with a
content-extraction library, and passes the title and text to \texttt{GPT-4o}
under the \texttt{NewsExtraction} schema to obtain the mentioned tickers, a
sentiment label, and a short Vietnamese summary (Section~\ref{sec:news}). Each
article is de-duplicated by its title and link and discarded if its body is too
short, then written to the store and linked to every mentioned stock through the
\texttt{Article\_Stock} table. As with price ingestion, requests are paced to
remain within provider rate limits.

\subsection{Safety and robustness}
Several mechanisms harden the system against both adversarial and accidental
failure. Every model-generated SQL statement is checked by a sanitizer that
admits only read-only \texttt{SELECT}s over an allow-list of tables and runs them
under a timeout, so the conversational assistant can never mutate the database
(Section~\ref{sec:chatbot}). Throughout the analysis pipeline, missing data is
signalled by explicit ``no data'' markers that the aggregator detects, so a
single absent source degrades gracefully instead of aborting the run. Finally,
the aggregator's decisive computations are performed in Python and override the
LLM, and any unhandled error resolves to a safe \textsf{Wait} response rather
than a malformed one.

\subsection{Mobile client}
The front end is a React Native application built with Expo, using file-based
routing. Authentication tokens are kept in encrypted device storage, and the
market screens render bespoke chart components---candlestick, line and bar
charts, and the RSI, MACD and KDJ panels---so that the indicators discussed in
Section~\ref{sec:technical} are presented to the user in the same vocabulary the
analysis pipeline uses.

\section{Evaluation}
\label{sec:evaluation}

We evaluate the analysis pipeline by backtesting its recommendations on the VN30
constituents and comparing them against two references: the unfiltered technical
engine that drives the pipeline, and a passive buy-and-hold strategy.

\subsection{Experimental setup}
The backtest covers the VN30 universe (29 constituents with sufficient history)
over a one-year window of daily candles, from January 2024 to January 2025,
with the VN-Index as the market benchmark. To avoid look-ahead bias, indicators
at bar $i$ use only data up to $i$ and the entry price is the open of the next
bar. To bound the (LLM) cost, the full pipeline is invoked only on dates where
the technical engine emits a signal; on those dates the aggregator produces the
buy/wait decision and the risk-managed targets of Section~\ref{sec:aggregator}.
Each realized trade contributes a percentage return, from which we compute
per-symbol performance, benchmark comparisons, and significance tests using the
framework in \texttt{backend/app/backtest}. Unless stated otherwise, aggregate
figures are reported as the mean across symbols, with the median given where the
distribution is heavily skewed by outliers.

\subsection{Aggregate performance}
Table~\ref{tab:perf} summarizes performance across the universe. The most
visible effect of the pipeline is \emph{selectivity}: by requiring a clear
technical trigger and vetoing on weak fundamentals or negative news, it takes
$134$ trades in total versus $224$ for the unfiltered technical engine---roughly
$40\%$ fewer. The surviving trades have a higher average return per trade
($4.3\%$ vs.\ $2.4\%$) and a markedly smaller average maximum drawdown ($-9.7\%$
vs.\ $-13.4\%$), consistent with the risk-averse design. The trade-off is that
the unfiltered engine, by participating more often, attains a higher aggregate
total return and Sharpe ratio over this particular sample.

\begin{table}[t]
\centering
\caption{Aggregate backtest performance over the VN30 universe (per-symbol mean
unless noted). ``Technical-only'' is the unfiltered signal engine; buy-and-hold
is passive. Returns are skew-prone, so totals are given as mean (median).}
\label{tab:perf}
\begin{tabular}{@{}lccc@{}}
\toprule
Metric & LLM pipeline & Technical-only & Buy-and-hold \\
\midrule
Total trades            & 134           & 224           & --- \\
Win rate                & 50.1\%        & 54.5\%        & --- \\
Avg.\ return per trade  & 4.3\%         & 2.4\%         & --- \\
Total return            & 15.2\% (8.1\%) & 27.6\% (9.9\%) & 56.9\% (18.7\%) \\
Sharpe ratio            & 0.31          & 0.61          & --- \\
Profit factor (median)  & 1.34          & 1.73          & --- \\
Max drawdown            & $-9.7\%$      & $-13.4\%$     & --- \\
\bottomrule
\end{tabular}
\end{table}

\subsection{Benchmark comparison}
Relative to the unfiltered engine, the pipeline improved total return on $15$ of
$29$ symbols (median impact $+0.4\%$), i.e.\ the LLM filter helped about as often
as it hurt while substantially reducing drawdown. Against buy-and-hold the
picture is less favourable: only $11$ of $29$ symbols beat the passive strategy
(median return gap $-7.5\%$). This is expected for a selective, partially
invested, long-only strategy evaluated over a predominantly bullish sample, in
which simply staying fully invested is hard to beat. Compared with random entry
strategies of the same trade count, the pipeline sat at the $61\%$ percentile on
average---better than chance, but not decisively so.

\subsection{Behaviour across market regimes}
Table~\ref{tab:regime} breaks trades down by market regime. Strikingly, the
pipeline opened \emph{no} positions during downtrends, concentrating its activity
in sideways ($83$ trades) and uptrending ($51$ trades) markets; it was most
effective in sideways regimes ($53.0\%$ win rate, $4.0\%$ average return). This
behaviour follows directly from the buy/wait logic, in which only clearly
positive evidence triggers an entry---an implicit, and desirable, avoidance of
falling markets.

\begin{table}[t]
\centering
\caption{Trades by market regime (pooled across symbols, trade-weighted). The
pipeline took no positions in downtrends.}
\label{tab:regime}
\begin{tabular}{@{}lccc@{}}
\toprule
Regime & Trades & Win rate & Avg.\ return per trade \\
\midrule
Uptrend   & 51 & 37.3\% & 1.7\% \\
Sideways  & 83 & 53.0\% & 4.0\% \\
Downtrend & 0  & ---    & --- \\
\bottomrule
\end{tabular}
\end{table}

\subsection{Statistical significance}
We caution against over-interpreting the point estimates above. With an average
of only $\approx 4.6$ trades per symbol, the per-symbol tests are statistically
underpowered: a one-sample $t$-test of returns against zero is significant
($p<0.05$) for none of the $29$ symbols, as is a sign-permutation test. The mean
information coefficient between the pipeline's confidence and forward returns is
weakly positive ($0.08$, positive for $13$ of $29$ symbols) but not significant.
We therefore present these results as evidence that the pipeline behaves as
designed---selective, drawdown-aware, and regime-avoiding---rather than as a
claim of a statistically established return edge; establishing the latter would
require a larger universe and a longer evaluation window
(Section~\ref{sec:discussion}).

% ============================================================
% TODO Section 7 — Discussion
% TODO Section 8 — Conclusion
% ============================================================

\begin{credits}
\subsubsection{\ackname} A bold run-in heading in small font size at the end of the paper is
used for general acknowledgments, for example: This study was funded
by X (grant number Y).

\subsubsection{\discintname}
It is now necessary to declare any competing interests or to specifically
state that the authors have no competing interests. Please place the
statement with a bold run-in heading in small font size beneath the
(optional) acknowledgments\footnote{If EquinOCS, our proceedings submission
system, is used, then the disclaimer can be provided directly in the system.},
for example: The authors have no competing interests to declare that are
relevant to the content of this article. Or: Author A has received research
grants from Company W. Author B has received a speaker honorarium from
Company X and owns stock in Company Y. Author C is a member of committee Z.
\end{credits}
%
% ---- Bibliography ----
%
% BibTeX users should specify bibliography style 'splncs04'.
% References will then be sorted and formatted in the correct style.
%
% \bibliographystyle{splncs04}
% \bibliography{mybibliography}
%
\begin{thebibliography}{8}
\bibitem{vsd2025}
Vietnam Securities Depository and Clearing Corporation (VSDC): Securities account
statistics. Reported in: Vietnam's stock market has more than 11 million trading
accounts. \url{https://www.vietnam.vn/en/thi-truong-chung-khoan-viet-nam-co-hon-11-trieu-tai-khoan-giao-dich}, last accessed 2026/06/29

\bibitem{vinacapital2024}
VinaCapital: Vietnam's Resilient Stock Market. Insights report (2024).
\url{https://vinacapital.com/wp-content/uploads/2024/07/VinaCapital-Insights-Vietnams-Resilient-Stock-Market.pdf}, last accessed 2026/06/29

\bibitem{ftse2025}
FTSE Russell (LSEG): FTSE Russell announces results of September 2025 semi-annual
country classification review --- reclassification of Vietnam from Frontier to
Secondary Emerging market (announced October 2025).
\url{https://www.lseg.com/en/media-centre/press-releases/ftse-russell/2025/ftse-russell-country-classification-september-2025}, last accessed 2026/06/29

% NOTE: Author lists for the arXiv/journal entries below should be verified
% against the original papers before final submission.
\bibitem{finbertlstm2022}
Halder, S.: FinBERT-LSTM: Deep learning based stock price prediction using news
sentiment analysis. arXiv:2211.07392 (2022)

\bibitem{finbert2023}
Financial sentiment analysis using FinBERT with application in predicting stock
movement. arXiv:2306.02136 (2023)

\bibitem{phobertstock}
Stock article title sentiment-based classification using PhoBERT. In: CEUR
Workshop Proceedings, vol. 3026 (2021).
\url{https://ceur-ws.org/Vol-3026/paper25.pdf}

\bibitem{vnnews2023}
Sentiments extracted from news and stock market reactions in Vietnam.
International Journal of Financial Studies \textbf{11}(3), 101 (2023)

\bibitem{bloomberggpt}
Wu, S., et al.: BloombergGPT: A large language model for finance.
arXiv:2303.17564 (2023)

\bibitem{fingpt}
Yang, H., Liu, X.-Y., Wang, C.D.: FinGPT: Open-source financial large language
models. arXiv:2306.06031 (2023)

\bibitem{tradingagents}
Xiao, Y., et al.: TradingAgents: Multi-agents LLM financial trading framework.
arXiv:2412.20138 (2024)

\bibitem{finmem}
Yu, Y., et al.: FinMem: A performance-enhanced LLM trading agent with layered
memory and character design. arXiv:2311.13743 (2023)

\bibitem{vndl2025}
A comprehensive comparison of deep learning models for forecasting developing
market indices: Evidence from the Vietnamese stock market. Procedia Computer
Science (2025)

\bibitem{vnlstm2021}
Predicting Vietnamese stock market using the variants of LSTM architecture. In:
Lecture Notes in Networks and Systems. Springer (2021).
\doi{10.1007/978-3-030-92942-8_11}

\bibitem{vnml2024}
Applying machine learning algorithms to predict the stock price trend in the
stock market --- the case of Vietnam. Humanities and Social Sciences
Communications \textbf{11} (2024)

\bibitem{text2sqlsurvey}
A survey of large language model-based generative AI for text-to-SQL:
Benchmarks, applications, use cases, and challenges. arXiv:2412.05208 (2024)

\bibitem{spider2}
Lei, F., et al.: Spider 2.0: Evaluating language models on real-world enterprise
text-to-SQL workflows. arXiv:2411.07763 (2024)
\end{thebibliography}
\end{document}
